\scanpage{74}
qui est évidemment stable par $\mathrm{GL}(V)$ donc par $\mathrm{GP}(V)$ -- donc se descend en une base $\omega$ de $\underline{\det}M$, d'où de $\det\widetilde{\mathfrak g}$. Il est vrai que dans l'iso (46) il y a une convention de signe à faire -- à défaut, il y a l'ambiguïté des signes (qu'on pouvait définir une \deleted{\unsure{telle}} $\omega$ canoniquement modulo cette seule ambiguïté était évident, car sur $\mathbb Z$ tout Module inversible a \unsure{exactt|exactement} deux bases, différant par le signe -- et l'invariance par \deleted{$\unsure{\mathrm{GP}1|\mathrm{GP}(1)}$} $G_0$ est claire. Mais $\omega^{\otimes2}$ ne dépend plus de cette convention, ce qui donne un sens dénué de tte ambiguïté à
\[
\underset{\substack{\text{discriminant}\\\text{divisé}}}{\underline{\delta'(q)}}=-\omega^{\otimes2}\tag*{(48)}
\]
d'où
\[
\delta(\varphi)=\delta(q)=-2\omega^{\otimes2}\tag*{(49)}
\]
% archive sheet number: 73
